% ============================================================
%  estilo.tex — preámbulo y diseño del curso
% ============================================================

\usepackage[a4paper,top=26mm,bottom=26mm,inner=28mm,outer=22mm]{geometry}
\usepackage[T1]{fontenc}
\usepackage[utf8]{inputenc}
\usepackage[english]{babel}
\usepackage{lmodern}
\usepackage{textcomp}
\usepackage{microtype}
% disable <<, >>, -- ligatures in monospaced fonts
\DisableLigatures[<,>,']{encoding = T1, family = tt*}
\usepackage{xcolor}
\usepackage{graphicx}
\usepackage{longtable}
\usepackage{array}
\usepackage{booktabs}
\usepackage{tabularx}
\usepackage{enumitem}
\usepackage{listings}
\usepackage{tikz}
\usepackage[framemethod=TikZ]{mdframed}
\usepackage{titlesec}
\usepackage{fancyhdr}
\usepackage{needspace}
\usepackage{amsmath}
\usepackage{amssymb}
\usepackage[hidelinks,bookmarks=true,bookmarksnumbered=true]{hyperref}

\usetikzlibrary{shapes.geometric,arrows.meta,positioning,calc,fit,backgrounds,decorations.pathreplacing}

% ------------------------------------------------------------
%  Paleta
% ------------------------------------------------------------
\definecolor{ink}{HTML}{16202C}
\definecolor{accent}{HTML}{1F5FA9}
\definecolor{accentlight}{HTML}{E8F0FA}
\definecolor{green}{HTML}{1E7A4B}
\definecolor{greenlight}{HTML}{E6F4EC}
\definecolor{amber}{HTML}{A86A00}
\definecolor{amberlight}{HTML}{FDF3E0}
\definecolor{red}{HTML}{A32B2B}
\definecolor{redlight}{HTML}{FBEBEB}
\definecolor{purple}{HTML}{5B3C88}
\definecolor{purplelight}{HTML}{F1ECF8}
\definecolor{gray}{HTML}{5C6773}
\definecolor{graylight}{HTML}{F2F4F6}
\definecolor{grayborder}{HTML}{D3D9E0}
\definecolor{teal}{HTML}{0F6E70}
\definecolor{teallight}{HTML}{E4F2F2}

\color{ink}

% ------------------------------------------------------------
%  Encabezados y pies
% ------------------------------------------------------------
\pagestyle{fancy}
\fancyhf{}
\renewcommand{\headrulewidth}{0pt}
\renewcommand{\footrulewidth}{0pt}
\fancyhead[LE]{\footnotesize\color{gray}\nouppercase{\leftmark}}
\fancyhead[RO]{\footnotesize\color{gray}\nouppercase{\rightmark}}
\fancyfoot[LE,RO]{\small\color{accent}\bfseries\thepage}
\fancyfoot[RE,LO]{\footnotesize\color{gray}Operating Systems from Scratch}

\fancypagestyle{plain}{%
  \fancyhf{}%
  \renewcommand{\headrulewidth}{0pt}%
  \fancyfoot[LE,RO]{\small\color{accent}\bfseries\thepage}%
  \fancyfoot[RE,LO]{\footnotesize\color{gray}Operating Systems from Scratch}%
}

\renewcommand{\chaptermark}[1]{\markboth{\thechapter.\ #1}{}}
\renewcommand{\sectionmark}[1]{\markright{\thesection.\ #1}}

% ------------------------------------------------------------
%  Títulos
% ------------------------------------------------------------
\titleformat{\part}[display]
  {\normalfont\sffamily\huge\bfseries\color{accent}\filcenter}
  {\vspace*{2cm}\normalsize\mdseries\color{gray}PART \thepart}
  {14pt}
  {\Huge\color{ink}}
  [\vspace{6pt}{\color{accent}\rule{0.45\textwidth}{1.2pt}}]

\titleformat{\chapter}[display]
  {\normalfont\sffamily\bfseries}
  {\raggedright\normalsize\color{accent}CHAPTER \thechapter}
  {6pt}
  {\raggedright\LARGE\color{ink}}
  [\vspace{4pt}{\color{accent}\rule{\textwidth}{1pt}}]
\titlespacing*{\chapter}{0pt}{6pt}{26pt}

\titleformat{\section}
  {\normalfont\sffamily\large\bfseries\color{accent}}{\thesection}{0.7em}{}
\titleformat{\subsection}
  {\normalfont\sffamily\normalsize\bfseries\color{ink}}{\thesubsection}{0.7em}{}
\titleformat{\subsubsection}
  {\normalfont\sffamily\small\bfseries\color{gray}}{}{0em}{}

\setcounter{tocdepth}{1}
\setcounter{secnumdepth}{2}

% ------------------------------------------------------------
%  Listas
% ------------------------------------------------------------
\setlist[itemize]{leftmargin=1.5em,itemsep=2pt,topsep=4pt}
\setlist[enumerate]{leftmargin=1.8em,itemsep=2pt,topsep=4pt}
\setlist[description]{leftmargin=1.2em,style=nextline,font=\bfseries\color{accent}}

% ------------------------------------------------------------
%  Código
% ------------------------------------------------------------
\lstdefinestyle{base}{
  basicstyle=\ttfamily\footnotesize\color{ink},
  keywordstyle=\color{accent}\bfseries,
  commentstyle=\color{green}\itshape,
  stringstyle=\color{red},
  numberstyle=\ttfamily\tiny\color{gray},
  numbers=left,
  numbersep=9pt,
  showstringspaces=false,
  breaklines=true,
  breakatwhitespace=false,
  breakindent=1.2em,
  tabsize=2,
  columns=fullflexible,
  keepspaces=true,
  frame=leftline,
  framerule=2pt,
  rulecolor=\color{grayborder},
  backgroundcolor=\color{graylight},
  framexleftmargin=6pt,
  framexrightmargin=2pt,
  framextopmargin=4pt,
  framexbottommargin=4pt,
  xleftmargin=10pt,
  aboveskip=10pt,
  belowskip=10pt,
  literate=%
    {á}{{\'a}}1 {é}{{\'e}}1 {í}{{\'i}}1 {ó}{{\'o}}1 {ú}{{\'u}}1
    {Á}{{\'A}}1 {É}{{\'E}}1 {Í}{{\'I}}1 {Ó}{{\'O}}1 {Ú}{{\'U}}1
    {ñ}{{\~n}}1 {Ñ}{{\~N}}1 {ü}{{\"u}}1 {¿}{{?`}}1 {¡}{{!`}}1
}

\lstdefinelanguage{nasm}{
  morekeywords={mov,movzx,movsx,add,sub,mul,imul,div,idiv,inc,dec,and,or,xor,not,neg,
    shl,shr,sal,sar,cmp,test,jmp,je,jne,jz,jnz,jc,jnc,ja,jae,jb,jbe,jg,jge,jl,jle,
    call,ret,iret,iretd,push,pop,pusha,popa,pushad,popad,pushfd,popfd,lea,nop,hlt,cli,sti,
    in,out,inb,outb,lgdt,lidt,ltr,invlpg,int,lodsb,stosb,rep,times,db,dw,dd,dq,equ,resb,resw,resd,
    global,extern,section,bits,org,section},
  morekeywords=[2]{eax,ebx,ecx,edx,esi,edi,ebp,esp,ax,bx,cx,dx,si,di,bp,sp,
    al,bl,cl,dl,ah,bh,ch,dh,cs,ds,es,fs,gs,ss,cr0,cr2,cr3,cr4,eflags},
  sensitive=false,
  morecomment=[l]{;},
  morestring=[b]",
  morestring=[b]'
}

\lstdefinestyle{c}{style=base,language=C,
  morekeywords={uint8_t,uint16_t,uint32_t,uint64_t,int8_t,int16_t,int32_t,int64_t,
    size_t,uintptr_t,intptr_t,bool,true,false,NULL,inline,asm,volatile,restrict}}
\lstdefinestyle{asm}{style=base,language=nasm,
  keywordstyle=\color{accent}\bfseries,
  keywordstyle=[2]\color{purple}}
\lstdefinestyle{sh}{style=base,language=bash,numbers=none,
  backgroundcolor=\color{ink!5},keywordstyle=\color{purple}\bfseries}
\lstdefinestyle{mk}{style=base,language=make,numbers=none}
\lstdefinestyle{txt}{style=base,numbers=none,keywordstyle=\color{ink},
  backgroundcolor=\color{white},frame=none,xleftmargin=10pt}

\lstset{style=c}

\newcommand{\filepath}[1]{\texttt{\small\color{purple}\detokenize{#1}}}

% Cabecera de filetag encima de un bloque de código
\newcommand{\filetag}[1]{%
  \needspace{4\baselineskip}%
  \vspace{8pt}\noindent
  \begin{tikzpicture}
    \node[anchor=west,inner xsep=7pt,inner ysep=3pt,fill=accent,text=white,
          font=\ttfamily\scriptsize,rounded corners=1.5pt] {\detokenize{#1}};
  \end{tikzpicture}\par\nobreak\vspace{-8pt}}

% ------------------------------------------------------------
%  Cajas
% ------------------------------------------------------------
\newcommand{\boxtitle}[2]{{\sffamily\bfseries\color{#1}#2}\par\vspace{2pt}}

\newmdenv[
  linecolor=accent,backgroundcolor=accentlight,linewidth=0pt,
  leftline=true,rightline=false,topline=false,bottomline=false,
  innerleftmargin=10pt,innerrightmargin=10pt,innertopmargin=8pt,innerbottommargin=8pt,
  skipabove=10pt,skipbelow=10pt,leftmargin=0pt,rightmargin=0pt,
  linewidth=3pt
]{boxconcept}

\newmdenv[
  linecolor=amber,backgroundcolor=amberlight,linewidth=3pt,
  leftline=true,rightline=false,topline=false,bottomline=false,
  innerleftmargin=10pt,innerrightmargin=10pt,innertopmargin=8pt,innerbottommargin=8pt,
  skipabove=10pt,skipbelow=10pt
]{boxwarn}

\newmdenv[
  linecolor=green,backgroundcolor=greenlight,linewidth=3pt,
  leftline=true,rightline=false,topline=false,bottomline=false,
  innerleftmargin=10pt,innerrightmargin=10pt,innertopmargin=8pt,innerbottommargin=8pt,
  skipabove=10pt,skipbelow=10pt
]{boxcheck}

\newmdenv[
  linecolor=purple,backgroundcolor=purplelight,linewidth=3pt,
  leftline=true,rightline=false,topline=false,bottomline=false,
  innerleftmargin=10pt,innerrightmargin=10pt,innertopmargin=8pt,innerbottommargin=8pt,
  skipabove=10pt,skipbelow=10pt
]{boxanalogy}

\newmdenv[
  linecolor=red,backgroundcolor=redlight,linewidth=3pt,
  leftline=true,rightline=false,topline=false,bottomline=false,
  innerleftmargin=10pt,innerrightmargin=10pt,innertopmargin=8pt,innerbottommargin=8pt,
  skipabove=10pt,skipbelow=10pt
]{boxdanger}

\newmdenv[
  linecolor=teal,backgroundcolor=teallight,linewidth=3pt,
  leftline=true,rightline=false,topline=false,bottomline=false,
  innerleftmargin=10pt,innerrightmargin=10pt,innertopmargin=8pt,innerbottommargin=8pt,
  skipabove=10pt,skipbelow=10pt
]{boxrecipe}

\newmdenv[
  linecolor=grayborder,backgroundcolor=graylight,linewidth=0.8pt,
  roundcorner=3pt,
  innerleftmargin=10pt,innerrightmargin=10pt,innertopmargin=8pt,innerbottommargin=8pt,
  skipabove=10pt,skipbelow=10pt
]{boxplain}

\newenvironment{concept}[1]{\begin{boxconcept}\boxtitle{accent}{#1}}{\end{boxconcept}}
\newenvironment{warn}[1]{\begin{boxwarn}\boxtitle{amber}{#1}}{\end{boxwarn}}
\newenvironment{tryit}[1][Try it yourself]{\begin{boxcheck}\boxtitle{green}{#1}}{\end{boxcheck}}
\newenvironment{analogy}[1][An analogy]{\begin{boxanalogy}\boxtitle{purple}{#1}}{\end{boxanalogy}}
\newenvironment{danger}[1]{\begin{boxdanger}\boxtitle{red}{#1}}{\end{boxdanger}}
% Runnable recipe: #1 = the task, #2 = folder under examples/
\newenvironment{recipe}[2]{%
  \begin{boxrecipe}%
  {\sffamily\bfseries\color{teal}In practice: #1}\par
  {\ttfamily\scriptsize\color{gray}examples/#2/\quad---\quad make run}\par\vspace{5pt}}
  {\end{boxrecipe}}

\newenvironment{summary}[1][In short]{\begin{boxplain}\boxtitle{gray}{#1}}{\end{boxplain}}

% Bloque de "goals" al inicio de cada capítulo
\newenvironment{goals}{%
  \begin{boxplain}\boxtitle{accent}{By the end of this chapter you will know}%
  \begin{itemize}[leftmargin=1.3em,itemsep=1pt,topsep=0pt]}
  {\end{itemize}\end{boxplain}\vspace{4pt}}

% Ejercicios
\newcounter{exercise}[chapter]
\newcommand{\exercise}[1]{%
  \refstepcounter{exercise}%
  \medskip\noindent{\sffamily\bfseries\color{accent}Exercise \thechapter.\theexercise.}\ #1\par\medskip}

% Atajos tipográficos
\newcommand{\hi}[1]{\textbf{\color{accent}#1}}
\newcommand{\term}[1]{\textbf{#1}}
\newcommand{\reg}[1]{\texttt{\color{purple}#1}}
\newcommand{\hexa}[1]{\texttt{0x#1}}
\newcommand{\cod}[1]{\texttt{\small #1}}

% Decorative section break
\newcommand{\sectionbreak}{\par\vspace{6pt}\centerline{\color{grayborder}\rule{0.3\textwidth}{0.6pt}}\vspace{6pt}\par}

\setlength{\parskip}{0.45em}
\setlength{\parindent}{0pt}
\linespread{1.06}

\setlength{\emergencystretch}{2.5em}
\sloppy

\raggedbottom
